\section{Related Work}

% CITATION STATUS: author-year mentions below are positioning anchors and remain
% NEEDS-VERIFICATION until source-checked against a frozen .bib. No empirical
% claim is made here.

\paragraph{Generative-evaluation metrics.}
Fr\'echet Inception Distance (Heusel et al., 2017) and Inception Score (Salimans
et al., 2016) remain the default summaries of generative quality, despite known
estimator bias: FID is biased at finite sample sizes and sensitive to the
extrapolation to infinite samples (Chong and Forsyth, 2020). Kernel Inception
Distance (Bi\'nkowski et al., 2018) replaces the Fr\'echet form with an unbiased
MMD estimate; CMMD (Jayasumana et al., 2024) argues for CLIP embeddings with an
unbiased MMD; and FD-DINOv2 (Stein et al., 2023) exposes feature-space artifacts
and the unfair treatment of diffusion models. CertGen is deliberately
\emph{not} a new metric: it is a decision layer that consumes these metrics. We
keep Fr\'echet-form scores (FID, FD-DINOv2) descriptive because they are
nonlinear functionals of feature covariances without a clean bounded-mean
streaming estimator, and place rigorous certificates only on bounded-kernel
MMD/CMMD streams.

\paragraph{Evaluation reliability and reproducibility.}
Reported numbers are fragile to preprocessing: aliased resizing and other
pipeline subtleties materially change FID (Parmar et al., 2022). This motivates
our preprocessing lock and the requirement that a reported point estimate be
reproduced within tolerance \emph{before} any certificate is trusted. A separate
line studies benchmark and ranking reliability, largely in NLP (e.g., Dror et
al., 2018, on significance testing); generative-vision leaderboards have received
comparatively little such scrutiny, which is the gap CertGen targets.

\paragraph{Anytime-valid inference and e-processes.}
Our statistical machinery is borrowed, and we say so plainly. Time-uniform
confidence sequences (Howard et al., 2021), betting-based confidence sequences
for bounded means (Waudby-Smith and Ramdas, 2023), and the broader theory of
e-values, test (super)martingales, and e-processes (Vovk and Wang, 2021; Ramdas
et al., 2023) provide validity under optional stopping. For many simultaneous
decisions we use e-BH (Wang and Ramdas, 2022), which controls the false
discovery rate under \emph{arbitrary} dependence -- the relevant regime here,
since every model pair is compared against one shared reference set.

\paragraph{Sequential and kernel two-sample testing (closest prior work).}
The most direct neighbors are sequential two-sample and kernel tests:
betting/e-value two-sample and independence testing (e.g., Shekhar and Ramdas,
2023; Podkopaev and Ramdas, 2023) instantiate exactly the kind of anytime-valid
kernel comparison we rely on. We claim no methodological novelty over this line.
CertGen's contribution is orthogonal and applied: (i) framing generative-model
\emph{leaderboard wins} as a set of simultaneous, optional-stopping-safe
decisions on the paired discrepancy $\Delta_{AB}=d(A,R)-d(B,R)$; (ii) an
e-BH-controlled \emph{audit} of how many reported wins are statistically decided
at the sample sizes actually used; and (iii) a reproducibility discipline
(provenance ledger, preprocessing lock, reproduction gate) that makes such an
audit defensible on released samples and feature caches. In short, prior work
gives the test; we ask -- and answer with real benchmarks -- whether the field's
comparisons survive it.
